\documentclass[../../../include/open-logic-section]{subfiles}

\begin{document}

\olfileid{sth}{story}{blv}

\olsection{ملحق في قانون فرغه الأساسي الخامس}

سبق في~\olref[rus]{sec} أن راسل أقام مفارقته اعتراضًا على النسق الذي
بسطه فرغه في~\emph{Grundgesetze}. ولم يصرح ذلك النسق بالاستيعاب
الساذج؛ فنبين هنا بإيجاز ما \emph{صرح به}، ووجه اتصاله بالاستيعاب
الساذج ومفارقة راسل.

ونسق فرغه \emph{من الرتبة الثانية}، وغرضه صوغ معنى
\emph{امتداد مفهوم}.\footnote{يحاول فرغه، على وجه الدقة، إضفاء الصبغة
  الصورية على مفهوم أعم هو «مدى قيم» دالة. وامتدادات المفاهيم حالة خاصة
  من ذلك المفهوم الأعم. للتفاصيل، انظر
  \citet[pp.\ 8--9]{Heck2012}.} وباصطلاح مستمد منه نكتب
$\fregeext{{\OLId{latin-ordinary}{x}}}{{\OLId{latin-looped}{F}}({\OLId{latin-ordinary}{x}})}$ لـ\emph{امتداد المفهوم~${\OLId{latin-looped}{F}}$}. وهذه العملية تأخذ
\emph{محمولًا}، «${\OLId{latin-looped}{F}}$»، وتخرجه \emph{حدًا} من الرتبة الأولى، هو
«$\fregeext{{\OLId{latin-ordinary}{x}}}{{\OLId{latin-looped}{F}}({\OLId{latin-ordinary}{x}})}$». وبها عرّف فرغه الانتماء \emph{التعريف} الآتي:
\[
	{\OLId{latin-ordinary}{a}} \in {\OLId{latin-ordinary}{b}} \liff_\text{df} \exists {\OLId{latin-looped}{G}}({\OLId{latin-ordinary}{b}} = \fregeext{{\OLId{latin-ordinary}{x}}}{{\OLId{latin-looped}{G}}({\OLId{latin-ordinary}{x}})} \land Ga)
\]
ومعناه تقريبًا أن~${\OLId{latin-ordinary}{a}} \in {\OLId{latin-ordinary}{b}}$ إذا وفقط إذا اندرج~${\OLId{latin-ordinary}{a}}$ تحت مفهوم امتداده
${\OLId{latin-ordinary}{b}}$. والمكمّم «$\exists {\OLId{latin-looped}{G}}$» هنا من الرتبة الثانية. وأخذ فرغه كذلك
بالمبدأ المسمّى \emph{القانون الأساسي الخامس}:
$$\fregeext{{\OLId{latin-ordinary}{x}}}{{\OLId{latin-looped}{F}}({\OLId{latin-ordinary}{x}})} = \fregeext{{\OLId{latin-ordinary}{x}}}{{\OLId{latin-looped}{G}}({\OLId{latin-ordinary}{x}})} \liff \forall {\OLId{latin-ordinary}{x}} (Fx \liff Gx)$$
أي إن المفهومين، على التقريب، متحدا الامتداد إذا وفقط إذا تساويا
امتدادًا؛ وكل من «${\OLId{latin-looped}{F}}$» و«${\OLId{latin-looped}{G}}$» واقع مرة أخرى موقع المحمول. ومن هذين
الأصلين ينتج مبدأ بسيط يصل الانتماء بتحقق الصفة:

\begin{lem}[في \emph{Grundgesetze}]\ollabel{lem:Fregeextensions}
$\forall {\OLId{latin-looped}{F}} \forall {\OLId{latin-ordinary}{a}}({\OLId{latin-ordinary}{a}} \in \fregeext{{\OLId{latin-ordinary}{x}}}{{\OLId{latin-looped}{F}}({\OLId{latin-ordinary}{x}})} \liff Fa)$
\end{lem}

\begin{proof}
ثبّت~${\OLId{latin-looped}{F}}$ و${\OLId{latin-ordinary}{a}}$. فعندئذ يكون~${\OLId{latin-ordinary}{a}} \in \fregeext{{\OLId{latin-ordinary}{x}}}{{\OLId{latin-looped}{F}}({\OLId{latin-ordinary}{x}})}$ إذا وفقط إذا كان
$\exists {\OLId{latin-looped}{G}}(\fregeext{{\OLId{latin-ordinary}{x}}}{{\OLId{latin-looped}{F}}({\OLId{latin-ordinary}{x}})} = \fregeext{{\OLId{latin-ordinary}{x}}}{{\OLId{latin-looped}{G}}({\OLId{latin-ordinary}{x}})} \land Ga)$
(بحسب تعريف الانتماء)، وهذا إذا وفقط إذا كان
$\exists {\OLId{latin-looped}{G}}(\forall {\OLId{latin-ordinary}{x}}(Fx \liff Gx) \land Ga)$
(بالقانون الأساسي الخامس)، وهو إذا وفقط إذا كان~$Fa$، بالمنطق الأولي
من الرتبة الثانية.
\end{proof}

ومن ذلك يكاد الاستيعاب الساذج ينتج بلا واسطة:

\begin{lem}[في \emph{Grundgesetze}.]
$\forall {\OLId{latin-looped}{F}} \exists {\OLId{latin-ordinary}{s}} \forall {\OLId{latin-ordinary}{a}} ({\OLId{latin-ordinary}{a}} \in {\OLId{latin-ordinary}{s}} \liff Fa)$
\end{lem}

\begin{proof}
ثبّت~${\OLId{latin-looped}{F}}$. تعطي~\olref{lem:Fregeextensions} حينئذ
$\forall {\OLId{latin-ordinary}{a}} ({\OLId{latin-ordinary}{a}} \in \fregeext{{\OLId{latin-ordinary}{x}}}{{\OLId{latin-looped}{F}}({\OLId{latin-ordinary}{x}})} \liff Fa)$؛ فيلزم عنها
$\exists {\OLId{latin-ordinary}{s}}\forall {\OLId{latin-ordinary}{a}}({\OLId{latin-ordinary}{a}} \in {\OLId{latin-ordinary}{s}} \liff Fa)$ بالتعميم الوجودي. ولما كانت~${\OLId{latin-looped}{F}}$
اعتباطية تم الحكم.
\end{proof}

ثم تحصل مفارقة راسل إذا أخذنا~${\OLId{latin-looped}{F}}$ على الوجه الذي تعينه
$\forall {\OLId{latin-ordinary}{x}}(Fx \liff {\OLId{latin-ordinary}{x}} \notin {\OLId{latin-ordinary}{x}})$.

\end{document}



